% !TeX encoding = utf8
% !TeX spellcheck = fr

\documentclass[french, a4paper, 12pt]{refrep}
\settextfraction{1}

\usepackage{xcolor}
\usepackage{booktabs}
\usepackage{longtable}
\usepackage{tabularx}
\usepackage{rotating}
\usepackage{multicol}
\usepackage{multirow}
\usepackage[disable]{todonotes}
\usepackage[colorlinks=true]{hyperref}
\usepackage{gensymb}
\makeatletter

% Reference colors
\hypersetup{linkcolor=blue,     % internal links
linktocpage=true,               % only page numbers
urlcolor=blue,                  % external links 
bookmarks=true,
bookmarksnumbered=true,
breaklinks=true                 % wrap links is Ok
}

% Some shortcuts
\newcommand\xc{\textsf{XCSoar }}   
\newcommand\fl{\textsf{Flarm }}
\newcommand\al{\textsf{Altair }}

% Define command to insert XCSoar website
\newcommand{\xcsoarwebsite}[1]{\url{http://www.xcsoar.org#1}}

% Define command to insert tip image
\newcommand{\tip}[0]{\marginlabel{\parbox{1.1cm}{\includegraphics[width=0.7cm]{figures/reminder.pdf}}}}

% Define command to insert gesture image
\newcommand{\gesture}[1]{\marginlabel{{\it#1
}\parbox{1.3cm}{\includegraphics[width=0.7cm]{figures/gesture.pdf}}}}

% Define command to insert specific gesture image
\newcommand{\gesturespec}[1]{\marginlabel{\parbox{1.3cm}{\includegraphics[width=0.9cm]{figures/#1.png}}}}

% Define command to insert warning image
\newcommand{\warning}[0]{\marginlabel{\parbox{1.3cm}{\includegraphics[width=0.9cm]{figures/warning.pdf}}}}

% Define command to insert Achtung image
\newcommand{\achtung}[0]{\marginlabel{\parbox{1.3cm}{\includegraphics[width=2.5em]{figures/warning.pdf}}}}

% Define command to insert a flash image
\newcommand{\blitz}[0]{\marginlabel{\parbox{1.3cm}{\includegraphics[height=2.0em]{figures/reminder.pdf}}}}

% Define command to insert a stop
\newcommand{\halt}[0]{\marginlabel{\parbox{1.3cm}{\includegraphics[height=2.0em]{figures/warning.pdf}}}}

% Define command to reference a configuration item
\newcommand{\config}[1]{\marginlabel{\ref{conf:#1}
\parbox{1.3cm}{\includegraphics[width=0.8cm]{figures/config.pdf}}}}

% Define command to draw a sketch on the margin
\newcommand{\sketch}[1]{\marginpar{\parbox{4.0cm}{\includegraphics[angle=0,width=1.0\linewidth,keepaspectratio='true']{#1}}}}
\newcommand{\smallsketch}[1]{\marginpar{\includegraphics[angle=0,keepaspectratio='true']{#1}}}


% Potentially overdue ``InfoBox'' style macro 
\newcommand{\InfoBox}[0]{{InfoBox}}

% Enumerated todo's for the todonotes package
\newcounter{todocounter}
\newcommand{\todonum}[2][]{\stepcounter{todocounter}\todo[#1]{\thetodocounter: #2}}

\maxipagerulefalse

% Include XCSoar header and footer settings
% Colors
\definecolor{buttongray}{rgb}{0.831,0.816,0.784}


% A set of boxes, buttons etc.
%
% Simple gray button
\newcommand{\blink}[0]{$\triangleright$}
\newcommand{\bmenug}[1]{
	\fcolorbox {black}{buttongray}{{\footnotesize\textsf{#1}}}
}
\newcommand{\button}{\bmenug}
\newcommand{\bmenu}{\bmenug}
\newcommand{\bmenuw}[1]{
    \fcolorbox {black}{white}{{\footnotesize\textsf{#1}}}
}
\newcounter{mboxwidth}
\setcounter{mboxwidth}{14}
\newcommand{\setbuttonwidth}[1]{\setcounter{mboxwidth}{#1}}
\newcommand{\bmenut}[2]{    % Menu button (two lined)
    \fcolorbox {black}{buttongray}{
    \makebox[\value{mboxwidth}mm][c]{
        \begin{tabular}{c}
        {\footnotesize\textsf{#1}}\\
        {\footnotesize\textsf{#2}}
        \end{tabular}
    }
  }
}
\newcommand{\bmenuth}[3]{    % Menu button (three lined)
    \fcolorbox {black}{buttongray}{
    \makebox[\value{mboxwidth}mm][c]{
        \begin{tabular}{c}
        {\footnotesize\textsf{#1}}\\
        {\footnotesize\textsf{#2}}\\
        {\footnotesize\textsf{#3}}
        \end{tabular}
    }
  }
}
\newcommand{\bmenus}[1]{    % Menu button (single line)
    \fcolorbox {black}{buttongray}{
    \makebox[\value{mboxwidth}mm][c]{
      \begin{tabular}{c}
        {\footnotesize\textsf{#1}}\\
          \\
      \end{tabular}
      }
    }
}
\newcommand{\infobox}[1]{    % Normal Info box in text
    \fcolorbox {black}{white}{\makebox[1.7cm][c]{\textsf\strut #1}}
}

% Some more convenience
%
\newenvironment{jspecs}{    % Description spacing
\itemsep=2pt\topsep=3pt\partopsep=3pt\parskip=0pt
\begin{description}
\itemsep=2pt\topsep=3pt\partopsep=3pt\parskip=0pt
}
{\end{description}}

\newcommand{\jindent}[2]{   % Extra list spacing
  \noindent\makebox[0pt][r]{{#1}\hspace*{\marginparsep}}
  \parbox[t]{0.95\linewidth}{#2}\par
}

\widowpenalty=1000
\clubpenalty=1000

% the command \version prints the XCSoar version number
\newcommand{\version}{\begingroup\catcode`\_=\active\input{VERSION.txt}\endgroup}

% Define command to put a menu label on the margin
% aligned left
\newcommand{\menulabel}[1]{\marginpar{\parbox{5.0cm}{\raggedright #1}}}
% aligned right
\newcommand{\menulabelr}[1]{\marginpar{\parbox{4.05cm}{\raggedleft #1}}}

% Define some colors
\definecolor{AirspaceYellow}{rgb}{.99,.99,.19}
\definecolor{AirspaceRed}{rgb}{.99,.19,.19}




\usepackage{wasysym}
\usepackage{graphicx}        % -> includegraphics
\usepackage{paralist}
\usepackage{framed}
\usepackage{babel}
\hypersetup{
    pdftitle={\xc{} --- Prise en main},
    pdfauthor={An Onymous},
    pdfsubject={Démarrage rapide},
    pdfproducer={xcsoar.org},
    pdfkeywords={XCSoar}{Démarrage rapide}{Vol à voile}{Arrivée},
}
\parindent0mm                %no indent
\textwidth=170mm             %match print and screen
\setlength{\baselineskip}{1em}

% Set the page title
\usepackage{calc}
\usepackage{fancyhdr}
\pagestyle{fancy}

% Add XCSoar User Manual title to the header
\settototalheight\headheight{\mycontent}
\fancyhead[L]{{\mycontent}}

% Add page number to the footer (centered)
\fancyfoot{}
\fancyfoot[R]{\thepage}

% No line between content and header
\renewcommand{\headrulewidth}{0pt}

\fancypagestyle{plain}{
  % Clear fancy header and footer
  \fancyhf{}

  % Add page number to the footer (centered)
  \fancyfoot[R]{\thepage}

  % No line between content and header
  \renewcommand{\headrulewidth}{0pt}
}

\xcsoarheader{\xc{} --- Prise en main}

%
%##############  Colour, Background, Border ##############
%
\fboxrule0.4mm               % width Fbox-border
\definecolor{flashyellow}{rgb}{.9,.85,0}
\definecolor{flashblue}{rgb}{0,0,1}
%
%##############  advanced Macros and Menus ..##############
%
\newcommand{\nav}[3]{\bmenut{Nav}{#1/2}{\LARGE$\triangleright$}~\bmenut{#2}{#3}}%NavMenu
\newcommand{\display}[3]{\bmenut{Display}{#1/2}{\LARGE$\triangleright$}~\bmenut{#2}{#3}}%DisplayMenu
\renewcommand{\config}[3]{\bmenut{Config}{#1/3}{\LARGE$\triangleright$}~\bmenut{#2}{#3}}%ConfigNavMenu
\newcommand{\info}[3]{\bmenut{Info}{#1/3}{\LARGE$\triangleright$}~\bmenut{#2}{#3}}%InfoMenu
%
%
%##############  Document start of flash ##############
%
\setcounter{secnumdepth}{1}
\renewcommand*{\thesection}{\textsf{\arabic{section}}}


\begin{document}
%
%print-testchart 
%\tiny
%    \definecolor{mauve}{RGB}{224 176 255}
%    \begin{testcolors}[rgb,RGB,cmyk,hsb,HSB,HTML,gray,Gray]
%    \testcolor{black}
%    \testcolor{white}
%    \testcolor{darkgray}
%    \testcolor{gray}
%    \testcolor{lightgray}
%    \testcolor{red}
%    \testcolor{green}
%    \testcolor{blue}
%    \testcolor{cyan}
%    \testcolor{magenta}
%    \testcolor{yellow}
%    \testcolor{brown}
%    \testcolor{lime}
%    \testcolor{olive}
%    \testcolor{orange}
%    \testcolor{pink}
%    \testcolor{purple}
%    \testcolor{teal}
%    \testcolor{violet}
%    \testcolor{mauve}
%    \testcolor[cmyk]{0 0.5 1 0.42}
%    \testcolor{flashblue}
%    \testcolor{flashyellow}
%    \end{testcolors}
%\newpage
%
%
%##############  Titlepage of flash ##############
\thispagestyle{empty}
\begin{center}
\fontsize{34}{0}
\selectfont\textbf{\xc{} --- Prise en main}
\fontsize{12}{12}
\vspace{0.1em}
\includegraphics[angle=0,width=0.58\linewidth,keepaspectratio='true']{figures/blitzlogo.png}

\vspace{1em}
{\Huge Préparez-vous à l'embarquement\\}
{\Large avec votre\\}
{\Huge Assistant Tactique Personnel\\pour le Vol à Voile\\}
\end{center}

%##############  XCSoar Checklist of flash ##############
\newpage
%%%%%--------------------------------------------------------------------%%%%%
%%%%%--------------------------------------------------------------------%%%%%
%%%%%--------------------------------------------------------------------%%%%%
%%%%%--------------------------------------------------------------------%%%%%
%%%%%--------------------------------------------------------------------%%%%%

\begin{framed}
\begin{center}
	\xc{} est entièrement développé par des bénévoles.\\
	Cette documentation aussi.\\
	Si vous voyez des imperfections, vous pouvez facilement les faire disparaître~:\\
	\xcsoarwebsite{/develop/}

\end{center}
\end{framed}

\begin{center}
{\Huge\textbf{Checklist d'\xc}}
\end{center}
\section*{Éclair \textcolor{flashyellow}{On}}\label{ch:flashon}
%%%%%--------------------------------------------------------------------%%%%%

\subsection{\textcolor{flashblue}{Mise en piste d'\xc}}
\begin{compactitem}
\item acquérir un appareil et y installer \xc
\item récupérer les fichiers de données adaptés à votre secteur de vol
\item dire à \xc{} quels fichiers utiliser
\item dire à \xc{} la polaire et la masse de votre planeur
\item éventuellement, se connecter aux instruments
\item terminer la configuration et se familiariser
\item monter l'appareil
\item ajouter les points listés ci-dessous dans vos checklists
\item définir le point de départ
\end{compactitem}

\subsection*{\textcolor{flashblue}{Faire les vérifications pré-vols, incluant~:}}
\begin{compactitem}
\item configurer la polaire et la masse
\item configurer le vent et les paramètres de vol (MC, moucherons, QNH)
\item éventuellement, configurer un circuit
\end{compactitem}

\subsection*{\textcolor{flashblue}{Faire les vérifications avant le décollage, incluant~:}}
\begin{compactitem}
\item vérifier le vent et la configuration de vol une dernière fois
\end{compactitem}

\subsection*{\textcolor{flashblue}{Voler, prendre du plaisir}}

\subsection*{\textcolor{flashblue}{Faire les vérification post-vol}}
\begin{compactitem}
\item Télécharger les enregistrements du vol depuis l'enregistreur, les mettre sur skylines, la netcoupe et l'OLC
\item Récupérer les données statistiques sur le vol.
\end{compactitem}

\section*{{\color[rgb]{.9,.85,0}Éclair} Off}

\hspace*{1cm} C'était trop bref~?\\
\hspace*{4cm} Juste une seconde...\\
\hspace*{6cm} Temps d'éclairement prolongé,\\
\hspace*{10cm} on continue~:


%##############  Start of main body of flash ##############
\setlength{\parskip}{0.3\baselineskip}
\newpage
%%%%%--------------------------------------------------------------------%%%%%
%%%%%--------------------------------------------------------------------%%%%%
%%%%%--------------------------------------------------------------------%%%%%
%%%%%--------------------------------------------------------------------%%%%%
%%%%%--------------------------------------------------------------------%%%%%
\section{Installation d'\xc}\label{ch:XCSinstall}
%%%%%--------------------------------------------------------------------%%%%%

%\begin{wrapfigure}[10]{R}{0.23\textwidth}
%\vspace{-5.5em}
%  \begin{flushright}
%    \includegraphics[width=0.17\textwidth]{figures/xy01.jpg}
%    XY, ready for installation
%  \end{flushright}
%\end{wrapfigure}

Identifiez la distribution d'\xc{} adaptée à votre appareil et
système d'exploitation \textsf{(Android, Windows, Linux)} 
sur \xcsoarwebsite{/download/}.
Si vous avez choisi un répertoire d'installation autre que celui par défaut, assurez-vous que les droits
d'accès sont adaptés. Créez un sous-répertoire \texttt{XCSoarData} et copiez-y tous
les fichiers de données, c'est-à-dire les fichiers de configuration, d'espaces aériens, de points de virage et
de relief XCM.

En utilisant les valeurs par défaut des installateurs pour Android, Windows Mobile et PC, cela devrait
fonctionner sans problème. Afin de garder les choses simples, restez avec les valeurs par défaut.
L'installation sur une carte SD est possible.

\begin{compactitem}
\item[1.] Télécharger la dernière version d'\xc{} publiée sur la page d'accueil 
\xcsoarwebsite{} ou via l'application Google Play pour Android.
\item[2.] Télécharger le fichier de relief \texttt{name.xcm} depuis
\xcsoarwebsite{/download/maps/}
correspondant à votre secteur de vol. Si un fichier \texttt{.xcm} semble
défectueux, utilisez un autre navigateur (il y a un problème avec plusieurs versions
d'internet explorer, qui désintègrent les fichiers \texttt{.xcm}).
\item[3.] Télécharger des points de virage dans l'un des formats connu d'\xc~:
\texttt{.dat, .cup, .wpz, .wpt}.  Par défaut, allez les chercher sur
\xcsoarwebsite{/download/data.html}.
\item[4.] Télécharger un fichier d'espaces aériens \texttt{.txt} et/ou \texttt{.sua}. Là aussi, suivez
le lien mentionné ci-dessus.
\end{compactitem}

\subsection{\textcolor{flashblue}{Lancer \xc{} pour la première fois}}
Lancer \xc. Lors du démarrage, \xc{} vous demande d'entrer soit dans la vie réelle (VOL) soit
dans le mode de simulation (SIM). En mode SIM, vous pouvez préciser manuellement la vitesse, le cap et la hauteur
afin de vous exercer au sol. La position initiale en mode SIM
peut être définie en utilisant le menu d'action Circuit \bmenuw{Définir comme nouveau point de départ}.
En mode VOL, les valeurs des instruments des appareils déjà connectés sont fournis à \xc.
Rappelez-vous en lorsque vous tentez de vous connecter aux instruments. En mode SIM, toutes
les connexions sont arrêtées.

Démarrez en mode VOL. Dites à \xc{} quels fichiers utiliser. Accédez à la configuration du système à travers le menu
principal \button{Config}.

\begin{flushleft}\hspace*{1cm}\config{2}{System}{}\blink\bmenuw{Fichiers}
\blink\bmenuw{Fichiers}\\\end{flushleft}
\index{Carte}\index{Waypoints}\index{Espaces aériens}

Sélectionnez les fichiers appropriés dans la fenêtre de dialogue. Si aucun fichier ne s'affiche, assurez-vous une nouvelle fois que les fichiers
sont bien dans le sous-répertoire \texttt{/XCSoarData}.

\subsection{\textcolor{flashblue}{Gestionnaire de fichiers}}\index{Gestionnaire de fichiers}
\textbf{Notez que} Le gestionnaire de fichier sur Android nécessite Android~2.2
ou supérieur~!!!

Si l'appareil que vous avez choisi est capable d'accéder directement à internet, vous pouvez
facilement télécharger et mettre à jour les fichiers de données~:
\begin{flushleft}\hspace*{1cm}\config{2}{Gestionnaire}{de fichiers}\\\end{flushleft}
Une liste des fichiers déjà téléchargés s'affiche. Lors de la toute première fois, la
liste peut apparaître vide. Choisissez \blink Ajouter.

Une énorme liste de fichiers s'affichent. Sinon, vérifier votre connexion internet. L'équipe
d'\xc{} met à jour cette liste de fichiers et les URL correspondantes
afin de vous donner une liste aussi complète que possible des téléchargements
dans le monde (notez que le nom d'un même fichier peut être différent si vous le téléchargez par un autre moyen).
Pour ajouter un fichier sélectionner \blink Ajouter. 
La sélection d'un fichier démarrera immédiatement le téléchargement. Pour ajouter d'autres fichiers, faites de nouveau
\blink Ajouter.
Mettre à jour des fichiers se fait avec \blink Télécharger, dans le même écran du menu. Pour
débuter, téléchargez un seul fichier de relief, d'espaces aériens et de point de virage. Avec 
certaines cartes, par exemple celle affichant des zones et des points d'ascendances, la configuration peut
prendre un peu de temps à comprendre. Pour plus d'informations référez-vous s'il vous plaît au manuel principal ou
à \url{http://www.glidinghotspots.eu/}.

\subsection{\textcolor{flashblue}{Langue (Polices, Sons)}}
Par défaut, XCSor reprend la langue du système d'exploitation. Si vous
en préférez une autre, il y en a une trentaine de disponible en septembre 2019 (désolé,
le Breton, le Basque et le Parisien sont absents...~:-)
\begin{flushleft}\hspace*{1cm}\config{2}{Système}{}\blink\bmenuw{Apparence}\blink
\bmenuw{Langue, Input}\\\end{flushleft}

Cette page du menu permet aussi de changer les polices et, après avoir préparé un
fichier de données particulier, de faire beaucoup plus en définissant des événements. Toutefois, ceci n'est pas
supposé faire partie d'un guide de démarrage rapide, c'était juste pour info.

%%%%%--------------------------------------------------------------------%%%%%
%%%%%--------------------------------------------------------------------%%%%%
%%%%%--------------------------------------------------------------------%%%%%
%%%%%--------------------------------------------------------------------%%%%%
%%%%%--------------------------------------------------------------------%%%%%


\section{Terminer la configuration}
%%%%%--------------------------------------------------------------------%%%%%

\subsection{\textcolor{flashblue}{Sélectionner une polaire}}

%\begin{wrapfigure}[7]{R}{0.22\textwidth}
%\vspace{-9em}
%  \begin{flushright}
%    \includegraphics[width=0.27\textwidth]{figures/xy02.png}
%    xy figuring out the missing colours of life without configuration.
%  \end{flushright}
%\end{wrapfigure}

\xc{} possède une longue liste de polaires prédéfinies. Pour en choisir une, l'éditer,
ou enregistrer votre polaire particulière, allez dans la configuration de l'aéronef en faisant
\begin{flushleft}\hspace*{1cm}\config{2}{Aéronef}{}\\\end{flushleft}
La liste des aéronefs de cette compilation provient la liste officielle de 2012 du Deutscher Aero Club.

\subsection{\textcolor{flashblue}{Connecter un enregistreur, un \fl, un vario, etc. Configuration NMEA}}

Pour fonctionner, \xc{} a besoin, à minima, de données en temps réel provenant d'un récepteur GPS. Vous devez dire à 
\xc{} lequel utiliser, même s'il est intégré à l'appareil. Pour s'y connecter, utilisez
\begin{flushleft}\hspace*{1cm}\config{2}{Périph.}{}\blink\bmenuw{Éditer}\\\end{flushleft}
De nombreux appareils portables possèdent ce type de compasant. Toutefois, c'est toujours
une bonne idée de se connecter à un ou plusieurs périphériques externes~:
\begin{compactitem}
\item un récepteur GPS spécialisé possède une bien meilleure réception, fournissant de bien
meilleures données pour les mesures et les calculs
\item un indicateur de vitesse fournit de manière exacte et rapide des estimations du vent
sans nécessiter de spiraler
\item un variomètre améliore l'assistant en thermique
\item un circuit peut être déclaré dans un enregistreur de vol IGC et, après l'atterrissage, l'enregistrement
du vol peut être téléchargé
\item certains variomètres permettent la synchronisation de la configuration du MacCready avec
\xc
\item un \fl{} (voire une entrée ADS-B) fournit des informations sur les autres
autour de vous --- alors que \xc{} n'effectue PAS de détection de collision.
\end{compactitem}

\subsection{\textcolor{flashblue}{NMEA}}
Vous vous demandez ce que c'est~? NMEA signifie ``National Marine Electronics Association'' (une association américaine de fabricants d'appareils électroniques maritimes), ayant
atteint le statut de norme mondiale. La norme définit les interfaces électriques
ainsi que les protocoles utilisés pour transmettre des données de navigation et
autres. Comme avec le mile marin, une fois de plus, les aviateurs adoptent des pratiques des marins.
Ainsi en va de leurs instruments\textellipsis

\subsection{\textcolor{flashblue}{GPS interne}}
Le récepteur GPS de votre propre appareil portable pourrait être de bonne qualité.
Pour différentes raisons, vérifiez avec attention si c'est vrai en vol~:

\begin{compactitem}
\item L'orientation en trois dimensions de votre appareil et de son antenne GPS
interne peuvent causer une mauvaise réception. Quand la position GPS est perdue, beaucoup
d'appareils portables dépendent fortement du aGPS (``assisted GPS''). Vous pourriez devoir attendre
plusieurs minutes pour une nouvelle position lors d'un vol.
\item aGPS est basé sur des informations transmises par un opérateur de réseau de téléphonie
mobile. Sa réception peut devenir mauvaise en l'air et ainsi affecter la précision de la
position GPS.
\item Des appareils bon marché ou mal configurés peuvent ``inventer'' des positions GPS. Des valeurs
dites ``interpolées'' et possiblement fausses peuvent perturber les algorithmes d'\xc.
\end{compactitem}

Si la réception GPS de votre appareil semble être fiable, félicitations! Vous
avez acheté un bon appareil. Vous n'avez pas besoin d'un périphérique en plus, vous possédez déjà
un système \xc{} de ``base''. Dans le cas contraire, la méthode préférée est de se connecter à un récepteur
GPS externe. Un bon choix est un enregistreur IGC. Un très bon choix est de se connecter à un \fl.
Avec un enregistreur IGC, vous pouvez déclarer des circuits et télécharger des enregistrements de vol
certifiés IGC, ce qui est nécessaire pour prendre part à l'OLC.
Se connecter à un \fl{} vous donnera des informations sur les membres de votre équipe,
sur le trafic et des données de hauteur barométrique --- ainsi que des enregistrements de vol certifiés IGC
si votre périphérique le permet.
Avec un \fl{} connecté, vous possédez un système \xc{} ``classique''. Intégrer d'autres
périphériques est un travail en soit et sort du cadre de ce guide de démarrage rapide.

\subsection{\textcolor{flashblue}{Attacher les choses ensembles}}
Actuellement, il y a trois façons de se connecter à d'autres instruments: utiliser soit
un cable série, soit un cable série-vers-USB, soit une connection sans fil par Bluetooth. Cette
dernière solution est la meilleure méthode pour les appareils Android. Elle nécessite cependant
ce que l'on appelle un ``Bluetooth Serial Port Profile (SPP)'' (``profil Bluetooth pour port série'').
Avec la version 4 d'Android et postérieure, Google a introduit de nombreux problèmes provenant du profil Bluetooth.
\xc{} étant gratuit, il n'y a presque aucun risque à tester les capacités de votre appareil. Au cas où cela fonctionne... ;-)

\textbf{Veuillez noter} qu'avant de connecter des appareils en Bluetooth, vous devez
apparier ces appareils. Les étapes requises dépendent de vos appareils.
Utiliser les logiciels de votre système d'exploitation. 

Pour utiliser des appareils équipés de ports USB, vous avez besoin d'un convertisseur série-vers-USB. Allez
sur le site d'\xc{} pour plus de détails~: \xcsoarwebsite{/hardware/}.

Si vous utilisez une connexion série, que les PDAs possèdent en général, faites 
attention: il y a trois types différents de connecteur physique communément utilisés.
Ils utilisent différents niveaux de signaux physiques~:
\begin{compactitem}
\item RS232 utilise un niveau \pm 12~Volt pleine gamme
\item TTL utilise une gamme de signaux en 5~Volt
\item LLL (``Low level logic'') utilise une gamme de signaux en 3~Volt
\end{compactitem}
D'un point de vue logique, ces connexions sont identiques et font exactement
la même chose.
Assurez-vous toutefois que les gammes des signaux des ports séries soient identiques. Dans le cas contraire, vous aurez besoin d'un convertisseur
de niveau. À nouveau, allez consulter le site internet d'\xc{} pour plus de détails.

Les choses sont plus simples si vous utilisez des configurations de périphériques recommandées. Un appareil très courant sous Android
est le Dell Streak~5. Sur la page consacrée aux matériels du site internet d'\xc{} vous trouverez la liste de plusieurs adaptateurs Bluetooth connus
pour fournir des connexions sans fil stables.

\subsection{\textcolor{flashblue}{Nom du pilote}}
\begin{flushleft}\hspace*{1cm}\config{2}{Système}{}\blink\bmenuw{Config.}\blink\bmenuw{Enregistreur}\\\end{flushleft}

\subsection{\textcolor{flashblue}{Paramètres d'horloge}}
\begin{flushleft}\hspace*{1cm}\config{2}{Système}{}\blink\bmenuw{Setup}\blink\bmenuw{Heure}\\\end{flushleft}

\subsection{\textcolor{flashblue}{Déplacer manuellement la carte}}
C'est le moment de considérer la carte. Vous pouvez la déplacer en entrant dans le mode panorama.
\begin{flushleft}\hspace*{1cm}\display{1}{Panor.}{On}\\\end{flushleft}
Selon votre appareil, utilisez ses boutons, cadrans, etc. --- ou même votre doigt --- pour
déplacer la carte.

\subsection{\textcolor{flashblue}{Définir le point de départ}}\index{Point de départ}
Quand \xc{} démarre, l'affichage de la carte est centré sur votre point de
''départ''. À moins que vous ayez défini un ``vrai'' circuit, la fonction route d'\xc{} vous
montrera la direction de votre point de départ. Si vous ne l'avez pas défini, \xc{} en créera un
automatiquement à partir de l'endroit détecté comme le point de décollage --- et
vous affichera la route du retour.
\begin{compactitem}
\item Déplacez la carte jusqu'au point désiré
\item Cliquez/tapotez sur l'écran (tactile) à proximité du point voulu.
\item La fenêtre ``Éléments de la carte à cet endroit'' s'affiche
\item Choisissez le point de virage.
\item Sélectionnez \bmenuw{Détails}.
\item Allez à la page suivante en touchant \blink
\item Sélectionnez \bmenuw{Définir comme nouveau point de départ}.
\end{compactitem}

Après en avoir terminé avec des paramétrages de base, vous êtes prêt pour
les vérifications pré-vol. À ce point, vous pourriez éprouver un fort intérêt
à vous familiariser avec tous les boutons et paramètres qu'\xc{} fournit.
Il se peut que toutes les fonctionnalités d'\xc{} vous fassent voir des X de partout~! Il
y a alors deux stratégies possibles pour résoudre cette situation~:
\begin{compactitem}
\item Ouvrez le manuel principal d'\xc{} et prenez votre temps.
\item Effectuez les dernières étapes ci-dessous. Puis rejouer un vol enregistré --- ou démarrer
\xc{} en mode simulateur --- afin de voir
l'écran d'\xc{} et se familiariser avec les affichages et les choses qui s'y produisent.
\end{compactitem}
La meilleure des stratégies étant de faire les deux\textellipsis

\subsection{\textcolor{flashblue}{Action -- regarder \xc{} rejouer un vol}}
Choisissez l'enregistrement d'un vol et laissez \xc{} remonter dans le temps. Copiez le fichier dans
le sous-répertoire \verb+/XCSoarData/logs+ et débutez la séquence.

\begin{flushleft}\hspace*{1cm}\config{2}{Rejouer}{}\blink\bmenuw{Fichier\textellipsis\ Démarrer}\\\end{flushleft}

À présent, vous êtes en mesure d'étudier au sol le comportement d'\xc. Alors que vous découvrez ceci,
vous pourriez ressentir le désir de modifier l'agencement de l'affichage et le contenu des
InfoBoxes. Il vous est suggéré de ne faire ce qui suit
que si vous avez un réel besoin de le faire. En effet, \xc{} est fourni avec une
pré-configuration vraiment bien pensée.

\subsection*{\textcolor{flashblue}{Agencement de l'affichage (portrait, paysage, etc.)}}
Pour modifier l'agencement de base de l'affichage~:
\begin{flushleft}\hspace*{1cm}\config{2}{Système}{}\blink\bmenuw{Apparence}\blink
\bmenuw{Agencement de l'écran}\\\end{flushleft}

\subsection*{\textcolor{flashblue}{Pour modifier le contenu d'une InfoBoxe}}
Pour modifier une série d'InfoBoxes et le contenu des InfoBoxes~:
\begin{flushleft}\hspace*{1cm}\config{2}{Système}{}\blink\bmenuw{Apparence}\blink
\bmenuw{Paramètres InfoBox}\\\end{flushleft}

\subsection{\textcolor{flashblue}{Fichiers de profil / Sauvegarde}}
Pour utiliser \xc{} dans différentes situations, environnements et planeurs, on peut 
créer des profils. Si \xc{} en trouve plus d'un au démarrage,
il vous demandera lequel utiliser. Activer un profil revient juste
à appuyer sur un bouton.
Ainsi, une fois que vous serez satisfait de tout ce que vous avez configuré, effectuez une sauvegarde de votre
fichier de profil. D'autres raisons d'utiliser plusieurs profils~: 
\begin{compactitem}
\item Vous ne possédez pas un planeur, mais vous volez avec tous les jouets de votre club. Afin de 
ne pas avoir, à chaque fois que vous changez de machine, à modifier la polaire et la masse du planeur,
et à reconfigurer les périphériques d'\xc{} avec les menus,
vous créez différents fichiers de profil.
\item Vous avez participé à un concours. Vous aviez alors défini un autre point de départ, un autre planeur, et
vous aviez reconfiguré le contenu de plusieurs pages d'écran. Maintenant, vous voulez revenir rapidement à
votre configuration de tous les jours.
\item Vous êtes un instructeur, et vous montrez de temps à autre votre \xc{} 
à vos élèves pour leur faire voir de quoi il en retourne. Ainsi, vous pourriez
avoir besoin de configurations bien plus simples que celle que vous utilisez habituellement.
\end{compactitem}
Pour sauvegarder des profils, faites deux copies de votre fichier de profil. 
Une première qui garde l'extension \verb+.prf+, une seconde où vous changez l'extension en
\verb+.bak+.
Pour repartir du dernier exemple, vous pourriez avoir~:
\begin{compactitem}
\item\verb+FormationEnASK13.prf+
\item\verb+FormationEnASK13.bak+
\item\verb+humphreyAuPaysDesMerveilles.prf+
\item\verb+humphreyAuPaysDesMerveilles.bak+
\end{compactitem}
Lorsqu'\xc{} charge un profil spécifique \verb+.prf+, ce fichier
va inclure tous les paramètres de configuration que vous avez sélectionné lorsque ce profil était actif.
Ainsi, un fichier de sauvegarde \verb+.bak+ est nécessaire pour préserver votre configuration de référence. Vous
en jugerez l'utilité par vous-même.

\subsection{\textcolor{flashblue}{Monter l'appareil}}
Monter l'appareil de votre choix peut s'avérer plus compliqué que
prévu. Dans tous les cas, le support de l'appareil doit être en mesure résister à des conditions très difficiles,
et ce que vous utilisiez un support spécial, un support à ventouse ou un truc vissé sur quelque chose.
Au cas où vous auriez l'idée d'essayer, tenir votre téléphone à la main est une très
mauvaise idée. Vous feriez alors mieux d'utiliser une bande Velcro. Une
façon rapide est grossière peut être de~:
\begin{compactitem}
\item acheter une protection de téléphone en silicone et y coller du Velcro au dos.
\item acheter une bande Velcro à vous la mettre autour de la cuisse, légèrement au-dessus du genou.
\item y placer votre téléphone (avec sa protection) juste avant de monter à bord.
\end{compactitem}
Ne tenez jamais à la main un appareil complexe pendant que vous pilotez~!

%%%%%--------------------------------------------------------------------%%%%%
%%%%%--------------------------------------------------------------------%%%%%
%%%%%--------------------------------------------------------------------%%%%%
%%%%%--------------------------------------------------------------------%%%%%
%%%%%--------------------------------------------------------------------%%%%%

\section{Vérification pré-vol}
%%%%%--------------------------------------------------------------------%%%%%

%\begin{wrapfigure}[4]{R}{0.5\textwidth}
%\vspace{-4.5em}
%  \begin{flushright}
%    \includegraphics[width=0.48\textwidth]{figures/xy03.jpg}
%    xy preflight check
%  \end{flushright}
%\end{wrapfigure}

À inclure à votre Checklist.

Si vous utilisez un seul appareil \xc{} dans plusieurs planeurs, commencez par configurer la
polaire de votre planeur. Puis configurez votre vol.
Faites soit~:

\subsection*{\textcolor{flashblue}{Configurer la polaire}}
\begin{flushleft}\hspace*{1cm}\config{2}{Aéronef}{}\hspace{1em}
\end{flushleft}

Soit~:

\subsection*{\textcolor{flashblue}{Choisir un profil}}
Redémarrez \xc{} et choisissez un autre profil ayant la bonne 
polaire et la bonne masse.

\subsection*{\textcolor{flashblue}{Configurer le vol}}
\begin{flushleft}\hspace*{1cm}\config{1}{Vol}{}\blink\bmenuw{Charge alaire / 
Moucherons / QNH / Temp max}\\\end{flushleft}

\subsection*{\textcolor{flashblue}{Configurer le vent}}
\begin{flushleft}\hspace*{1cm}\config{1}{Vent}{}\blink\bmenuw{Vent auto\textellipsis\ / 
Vitesse / Direction}\\\end{flushleft}

\subsection*{\textcolor{flashblue}{Définir un circuit~?}}
Après avoir défini votre point de départ, \xc{} vous indiquera toujours comment
y revenir, à moins que vous n'ayez créé un circuit. Comme vous pourriez vous y attendre, \xc
va ensuite vous guider sur votre circuit.
Pour créer et utiliser des circuits, une brève introduction est donnée à la fin de ce
document. Vous pouvez configurer soit un circuit ``simple'' soit un circuit ``réel''.


%%%%%--------------------------------------------------------------------%%%%%
%%%%%--------------------------------------------------------------------%%%%%
%%%%%--------------------------------------------------------------------%%%%%
%%%%%--------------------------------------------------------------------%%%%%
%%%%%--------------------------------------------------------------------%%%%%

\section{Vérification avant le décollage}
%%%%%--------------------------------------------------------------------%%%%%

%\begin{wrapfigure}[4]{R}{0.5\textwidth}
%\vspace{-4.5em}
%  \begin{flushright}
%    \includegraphics[width=0.48\textwidth]{figures/xy04.jpg}
%    xy checking wind
%  \end{flushright}
%\end{wrapfigure}

À inclure à votre Checklist.

\subsection{\textcolor{flashblue}{Vérifier le vol}}
\begin{flushleft}\hspace*{1cm}\config{1}{Vol}{}
%\blink\bmenuw{Wing loading / Bugs / QNH / Max temp}
\\\end{flushleft}

\subsection{\textcolor{flashblue}{Vérifier le vent}}
\begin{flushleft}\hspace*{1cm}\config{1}{Vent}{}
%\blink\bmenuw{Auto wind\dots / Speed / Direction}
\\\end{flushleft}


\section{Voler}
Au cours du vol, il y a deux écrans d'informations qui pourraient
vous intéresser.
\subsection*{\textcolor{flashblue}{État du vol}}
\begin{flushleft}\hspace*{1cm}\info{2}{État}{}\blink\bmenuw{Temps\textellipsis}\\\end{flushleft}
\subsection*{\textcolor{flashblue}{Analyse du vol}}
\begin{flushleft}\hspace*{1cm}\info{1}{Analyse}{}\bmenuw{\blink}\textellipsis\\\end{flushleft}

%%%%%--------------------------------------------------------------------%%%%%
%%%%%--------------------------------------------------------------------%%%%%
%%%%%--------------------------------------------------------------------%%%%%
%%%%%--------------------------------------------------------------------%%%%%
%%%%%--------------------------------------------------------------------%%%%%

\section{Vérification post-vol}
%%%%%--------------------------------------------------------------------%%%%%

\subsection{\textcolor{flashblue}{Télécharger l'enregistrement du vol}}
Télécharger les enregistrements de vol depuis l'enregistreur \emph{connecté} de votre choix.
\begin{flushleft}\hspace*{1cm}\config{2}{Périph.}{}\\\end{flushleft}
\begin{compactitem}
\item sélectionnez le port
\item appuyez sur \bmenuw{Télécharger le vol}
\item une liste des enregistrements de vol stockés sur votre enregistreur externe s'affiche.
Choisissez-en un
\item appuyez sur \bmenuw{Sélection}
\item attendez\textellipsis\ Le transfert se fait par une connexion série un peu ancienne
\item choisissez-en un autre ou appuyez sur \bmenuw{Non} pour terminer
\item selon votre enregistreur, vous devrez peut-être patienter un peu, une fois de plus. L'enregistreur redémarrera peut-être
et prendra son temps pour revenir en fonctionnement normal.
\end{compactitem}

%\begin{wrapfigure}[8]{R}{0.4\textwidth}
%\vspace{-3em}
%  \begin{flushright}
%    \includegraphics[width=0.39\textwidth]{figures/xy05.jpg}
%    xy analyzing
%  \end{flushright}
%\end{wrapfigure}

\textbf{Notez} qu'en cas de difficultés, la première chose à faire est de s'assurer que le pilote
que vous utilisez assure le téléchargement à partir de votre enregistreur. Au cas où vous utilisiez un système
en cascade (par ex. \fl-Vario-\xc), il pourrait être nécessaire de relier directement
le variomètre et de choisir un autre pilote dans \xc{} afin qu'il puisse accéder aux
téléchargements depuis le \fl{} vers \xc.

\subsection{\textcolor{flashblue}{Analyse du vol}}
\begin{flushleft}\hspace*{1cm}\info{1}{Analyse}{}\\\end{flushleft}
Utilisez \bmenuw{\blink} pour naviguer entre les écrans successifs affichant les données statistiques
acquises au cours du vol, comme l'estimation du résultat OLC, un barographe du vol\textellipsis\
 et même une analyse de la polaire si vous disposez d'un instrument
fournissant la vitesse air vrai.

\subsection{\textcolor{flashblue}{État du vol}}
\begin{flushleft}\hspace*{1cm}\info{2}{État}{}\blink\bmenuw{Temps\textellipsis}\\\end{flushleft}
Pour avoir accès à plus d'informations statistiques.

\subsection{\textcolor{flashblue}{Mettre en ligne un vol}}

Il y a trois portails qui attendent de voir les enregistrements de vos vols.
Skylines (\url{https://skylines.aero}), NetCoupe (\url{http://www.netcoupe.net})
et OLC (\url{http://www.onlinecontest.org}).
Si vous voulez y mettre en ligne un vol, probablement avez-vous intérêt à le faire immédiatement après le vol. L'OLC vous donne un délai de 24~heures pour y copier vos enregistrements ;  la NetCoupe deux semaines.
Skylines est un peu moins restrictif, et est votre portail pour le suivi en temps réel et\textellipsis\
vos fichiers de vols sont dans un sous-répertoire \texttt{logs}.

\subsection*{\textcolor{flashblue}{Tout éteindre}}
Afin d'économiser du jus, assurez-vous que vous avez bien pensé à \emph{tout} couper.

%%%%%--------------------------------------------------------------------%%%%%
%%%%%--------------------------------------------------------------------%%%%%
%%%%%--------------------------------------------------------------------%%%%%
%%%%%--------------------------------------------------------------------%%%%%
%%%%%--------------------------------------------------------------------%%%%%
\section{Circuits}
%%%%%--------------------------------------------------------------------%%%%%

\subsection{\textcolor{flashblue}{Fonction ``Retour au point de départ'', Fonction ``Aller à''}}
Après que vous ayez défini un ``vrai'' circuit pour la première fois, la fonctionnalité de base
de ``Retour au point de départ'' est définitivement supprimée. À la place, utilisez la fonction
``Aller à'', qui est le circuit le plus simple. \textbf{Notez qu'}il y a plusieurs fonctions, sous
forme d'InfoBoxe, liées au point de départ. Elles ne sont pas modifiées.
La seule fonction annulée est une \emph{astuce} qu'\xc{} utilise quand
il ne trouve pas de `vrai'' circuit défini par l'utilisateur --- ou jusqu'à ce que l'utilisateur définisse 
vraiment son tout premier circuit.

\subsection{\textcolor{flashblue}{Circuit simple~: ``Aller à''}}
Il y a plusieurs façons de créer un circuit simple. Dans tous les cas, sélectionnez un
point de virage et appuyez sur \bmenuw{Aller à}.

Cas les plus courants~:
\begin{flushleft}\hspace*{1cm}\nav{1}{Waypoint}{Liste}\\\end{flushleft}
\begin{compactitem}
\item Utilisez le menu de filtrage pour identifier le point de virage. Attention, vous pouvez paramétrer
plusieurs filtres. Si aucun point de virage n'apparaît, la combinaison des multiples filtres
pourrait vous empêcher de trouver votre point. 
\item Appuyez sur \bmenuw{OK}.
\item Appuyez sur \bmenuw{Sélection}, puis sur \bmenuw{Aller à}.
\end{compactitem}
Ou utilisez la fonction panorama de la carte, puis travaillez de manière similaire au cas ``Retour au point de départ''~:
\begin{compactitem}
\item Déplacez la carte jusqu'au point désiré
\item Cliquez/tapotez sur l'écran (tactile) à proximité du point voulu
\item La fenêtre ``Éléments de la carte à cet endroit'' s'affiche
\item Choisissez le point de virage
\item Appuyez sur \bmenuw{Aller à}
\end{compactitem}
 
\subsection{\textcolor{flashblue}{Circuit~: créer, sauver, charger, déclarer}}
Créer un nouveau circuit~:
\begin{flushleft}\hspace*{1cm}\nav{1}{Circuit}{}\blink\bmenuw{Gérer}\blink\bmenuw{Nouveau circuit}\\\end{flushleft}
\begin{compactitem}
\item Dans le même menu à l'écran, vous trouverez \bmenuw{Circuits enregistrés}, \bmenuw{Enregistrer} et 
\bmenuw{Déclarer} (à l'enregistreur). Après avoir appuyé sur \bmenuw{Nouveau circuit}, on vous demande de
confirmer, puis le menu concernant les \bmenuw{Règles} s'affiche. Définissez alors 
le \bmenuw{Type de circuit} en premier.
\item En sélectionnant \bmenuw{Point de virage} et \bmenuw{Ajouter point de virage} un écran de 
choix \bmenuw{Sélection Waypoint} s'affiche. Modifiez vos points de virage. En sélectionnant
\bmenuw{Éditer le point}, vous pourrez faire vos changements.
\item Ne fermez pas encore la fenêtre~! Pour terminer, utilisez  \bmenuw{Gérer} puis \bmenuw{Enregistrer},
et enfin \bmenuw{Fermer}.
\end{compactitem}

\subsection*{\textcolor{flashblue}{Circuit~: Heure de début et heure de fin}}
Comme vous arrivez à la fin de ce guide, pour creuser ce sujet, consulter le manuel principal d'\xc.

Merci pour le temps passé à lire ce guide et\textellipsis
%\begin{wrapfigure}[4]{l}{0.5\textwidth}
%\vspace{-1.3em}
%  \begin{centering}
%    \includegraphics[width=0.48\textwidth]{figures/xy06.jpg}
%  \end{centering}
%\end{wrapfigure}

\textellipsis\ faites de bons vols~!%no section at all, formatting reasons

%##############  Indexes of flash ##############
\begingroup
\fontfamily{ptm}
\normalsize
\fontseries{c}\selectfont
\setlength{\parskip}{0.05\baselineskip}
\tableofcontents
\endgroup



%##############  Closing page of flash ##############
\newpage
\pagestyle{empty}

\normalsize{Pour plus d'informations, consultez~:}

\begin{center}
    \includegraphics[angle=0,width=0.5\textwidth,keepaspectratio='true']{graphics/logo.png}
    \vskip 0.5cm
%   option 1: web release 
%    \includegraphics[angle=0,width=0.66\textwidth,keepaspectratio='true']{graphics/title.pdf}\\
%   option 2: print release
    \fontsize{50}{0}
    \selectfont\textbf{\xc}\\
    \fontsize{12}{12}
    \vspace{0.2em}
    \LARGE{La Navigation Open Source}\\
    \vspace{1.2em}
    \LARGE{Manuel de l'utilisateur}\\
    
\vspace{9em}

\end{center}
\begin{flushright}
\normalsize disponible sur~: \xcsoarwebsite{/discover/manual.html}
\end{flushright}

\end{document}
